% =====================================================================
% VERSION ENSAE, REDUCTION SFE : chapitre 2, revue de littérature et positionnement.
% Phase 3 : copie a texte constant de chapitres/03_etat_art_positionnement.tex.
% Phase 4 (1er octobre 2026) : reecrit et condense ; toutes les citations, le tableau
% comparatif et les nombres sont conserves. Le texte d'origine reste dans chapitres/
% (version Institut) et dans l'historique Git.
% Lu par main_ensae.tex seul : ne jamais le faire lire par main_institut.tex,
% qui garde chapitres/.
% =====================================================================
\chapter{Revue de littérature et positionnement}
\label{ch:positionnement}

% SOURCES-SCRIPTS: 08g 32 01 23

Les travaux mobilisés étudient l'assurabilité et la rareté des données
\citep{BienerElingWirfs2015}, les pertes extrêmes, la contagion entre entités et les arrivées dans
le temps. Nous reprenons leurs lois et leurs états successifs pour étudier la propagation entre
les domaines de contrôle d'une même entité.

\section{Comparaison aux approches existantes}
% SOURCES-SCRIPTS: 32 23 01

Le tableau~\ref{tab:positionnement-references} distingue ces objets. Notre articulation entre
contrôles d'une même entité et incertitude sur la direction n'est pas traitée conjointement dans
les références comparées. Les chapitres suivants la construisent à partir de ces briques.

\begin{table}[H]\centering
\begingroup
\definecolor{reperecitation}{HTML}{A52D39}
\hypersetup{colorlinks=true,citecolor=reperecitation}
\small
\newcommand{\objetcompare}[1]{\shortstack{\oui\\{\scriptsize\itshape #1}}}
\renewcommand{\arraystretch}{1.3}
\setlength{\tabcolsep}{3pt}
\begin{tabularx}{\textwidth}{@{}p{.29\textwidth}*{5}{>{\centering\arraybackslash}X}@{}}
\toprule
 & \multicolumn{2}{c}{\textbf{Modélisation du risque}} & \multicolumn{3}{c}{\textbf{Structure et décision}} \\
\cmidrule(lr){2-3}\cmidrule(lr){4-6}
\textbf{Références} & \shortstack{Lois de\\pertes} & \shortstack{Dynamique\\temporelle} & \shortstack{Piliers d'une\\même entité} & \shortstack{Identification\\partielle} & Remédiation \\
\midrule
\multicolumn{6}{@{}l}{\textit{Sévérité et dépendance cyber}} \\
\textcolor{reperecitation}{\citet{BohmeKataria2006,HerathHerath2011}} & \objetcompare{Copules} & & & & \\
\addlinespace
\textcolor{reperecitation}{\citet{elingwirfs2019,FarkasLopezThomas2021}} & \objetcompare{GPD} & & & & \\
\addlinespace
\textcolor{reperecitation}{\citet{HillairetLopez2021,HillairetLopezDOultremontSpoorenberg2022}} & & \objetcompare{États SIR} & & & \\
\addlinespace
\textcolor{reperecitation}{\citet{BessyRolandBoumezouedHillairet2021,BoumezouedCherkaouiHillairet2023}} & & \objetcompare{Hawkes} & & & \\
\midrule
\multicolumn{6}{@{}l}{\textit{Information incomplète et quantité d'intérêt}} \\
\textcolor{reperecitation}{\citet{Manski1990,Tamer2010}} & & & & \objetcompare{Ensembles} & \\
\midrule
\textbf{Ce travail} & \objetcompare{GPD reprise} & \objetcompare{États posés} & \objetcompare{Cascade} & \objetcompare{Exploration} & \objetcompare{Scénarios} \\
\bottomrule
\end{tabularx}
\endgroup
\caption{Objets traités par les références comparées. Une case vide ne signifie pas une
impossibilité générale de la méthode. La dernière ligne rassemble les briques reprises et leur
adaptation aux piliers DORA.}
\label{tab:positionnement-references}
\end{table}


Les matrices de risque et l'analyse des modes de défaillance, de leurs effets et de leur criticité
(AMDEC) classent les risques. L'indice RPN (\emph{Risk Priority Number}) multiplie gravité, occurrence et détection. Une copule relie les lois individuelles, appelées marges, pour former une loi jointe de pertes. Ni ce classement ni cette loi des seules pertes
ne décrivent l'ordre d'atteinte des piliers.

Les modèles de réseau représentent des transmissions entre entités, sous une structure et des
taux d'infection qui ne sont pas observés ici. Nous en reprenons l'idée d'états successifs,
appliquée aux contrôles d'une même entité. Les processus de Hawkes répondent à une autre question :
un incident augmente-t-il le taux d'arrivée d'autres incidents ? Les réseaux bayésiens décrivent
des liens dirigés sur un graphe acyclique avec des probabilités conditionnelles.
La section~\ref{sec:travaux-par-theme} détaille ces références ; le chapitre~\ref{chap:resultats} compare les effets de
modèles de dépendance sur le capital et le chapitre~\ref{ch:robustesse} examine l'information temporelle disponible.

\section{Dépendance dirigée et dépendance à l'ordre}
% SOURCES-SCRIPTS: 08g 32 01 23

L'objet propre à la cascade peut maintenant être précisé. La matrice $W$ de contagion entre
piliers, posée à dire d'expert puis bornée au chapitre~\ref{chap:partielle}, donne par $W_{jk}$
la propension d'une défaillance de $j$ à entraîner celle de $k$. La ligne désigne l'émetteur,
la colonne le récepteur. Elle est dirigée, au sens où $W_{jk}\neq W_{kj}$, et autorise des cycles
comme P1$\to$P2$\to\cdots\to$P1. La normalisation de Leontief les stabilise
($\rho(W)\le g<1$, forme réduite $(I-W)^{-1}$).

Avec le score de criticité du modèle, de $1$ à $10$ et construit à partir des barèmes d'expert
(chapitre~\ref{chap:cascade}), $22$ des $26$ sous-ensembles d'au moins deux piliers changent de
score selon l'ordre. L'écart atteint $6$ points : $8$ pour P1$\to$P2, contre $2$ pour
P2$\to$P1. Un score statique invariant par permutation ne peut pas reproduire cette propriété.
La non-transitivité de la dominance, testée sous trois formulations, ne se vérifie pas et n'est
pas revendiquée (chapitre~\ref{chap:cascade}).

\section{Travaux par thème}
\label{sec:travaux-par-theme}

\subsection*{Dépendance et contagion}
% SOURCES-SCRIPTS: 08g
La corrélation entre risques cyber a d'abord été traitée par
copules \citep{BohmeKataria2006, HerathHerath2011}. Sur la chronologie PRC, l'indice de Gini de la charge journalière vaut $0{,}69$ ($2\,922$ jours, $4\,614$ événements) : les sinistres se concentrent sur peu de jours, ce qui est compatible avec une accumulation par cause commune.
\citet{HillairetLopez2021} couplent des processus de comptage à un modèle épidémiologique
compartimental de type SIR (susceptibles, infectés, rétablis), et \citet{HillairetLopezDOultremontSpoorenberg2022} y ajoutent une
structure de réseau. \citet{XuHua2017} et \citet{Baldwin2017} modélisent la contagion entre cibles.

\subsection*{Dynamique temporelle}
\citet{BessyRolandBoumezouedHillairet2021} ajustent des processus de
Hawkes aux données PRC, l'une de nos sources, puis \citet{BoumezouedCherkaouiHillairet2023}
ajoutent un terme exogène. \citet{HillairetReveillacRosenbaum2021} en développent la théorie pour le
pricing de dérivés, et \citet{Peng2016} les appliquent à des taux d'attaque horodatés.

\subsection*{Fréquence, sévérité et queue}
\citet{Edwards2016} ajustent une binomiale négative au
nombre journalier de brèches, et \citet{LuZhangZhu2024} séparent fréquence et nombre de brèches par
événement, une structure compatible avec la loi de fréquence retenue ici (binomiale négative, chapitre~\ref{chap:socle}). \citet{FarkasLopezThomas2021}
tarifent par arbres de régression GPD. La queue lourde des pertes cyber est documentée
\citep{MaillartSornette2010, wheatley2016}, et \citet{elingloperfido2017} comparent plusieurs
familles de lois sur les brèches de données. \citet{elingwirfs2019} étudient les pertes cyber de la
base SAS OpRisk Global par les valeurs extrêmes, et leur étude est le précédent le plus proche pour la sévérité.

\subsection*{Travaux sur les mêmes sources}
\citet{Ranguin2024} modélise le risque cyber des
entreprises françaises pour la Caisse Centrale de Réassurance, à partir de PRC, VCDB, Hackmageddon et
de l'étude LUCY. La sévérité y est construite par arbres de régression GPD sur le volume de données
compromises. La sinistralité d'un portefeuille fictif est simulée, sinistres courants compris. Les lois
sont ensuite reparamétrées dans un cadre bayésien sur le marché français. Cette étude ne produit pas de
mesure de capital à un quantile réglementaire. L'objet étudié ici est le besoin de capital d'une
entité financière selon son état de conformité, avec une sévérité calibrée sur des pertes en montants
et une dépendance modélisée entre piliers.

Sur PRC, \citet{KherLopezRapior2023} estiment la gravité d'un incident à partir de sa description
textuelle par réseaux de neurones, et \citet{Dountio2026} prolonge cette démarche. Ces travaux
classent la sévérité d'un incident en volume de dossiers compromis.

\subsection*{Cadre structurel et identification partielle}
Le facteur latent commun vient du modèle à facteur de \citet{Vasicek2002}, et la normalisation qui stabilise une contagion avec
cycles provient de l'analyse entrées-sorties de \citet{Leontief1936}.

L'identification partielle caractérise des valeurs admissibles puis leurs conséquences sur une
quantité d'intérêt \citep{Manski1990,Manski1995,Tamer2010}. Ce principe est repris pour
explorer le capital à direction inconnue, conditionnellement à la force des liens posée.
Le chapitre~\ref{chap:partielle} précise l'ensemble admissible et l'échelle de contraintes de structure.

\subsection*{Cascade de risques climatiques}
Le réseau de cascade climatique de \citet{ccrn2026} sépare le temps calendaire des pas de
propagation. Cette architecture, le protocole de comparaison et l'expérience d'ablation en sont
repris. Ses valeurs synthétiques n'entrent dans aucune calibration de ce travail.

\section{Positionnement}

Le tableau~\ref{tab:apports-institut} relie chaque question à une méthode et au résultat qui
permet d'en juger. Il sert de guide aux chapitres suivants : construire le calcul, l'éprouver
quand l'information manque, puis comparer les choix de remédiation.

% HARNAIS-HORS-SECTION: guide de lecture contenant des renvois de chapitres et de sections,
% sans nouvelle valeur empirique ni nouveau parametre.
\begin{table}[H]\centering\small
\begingroup
\definecolor{reperecritique}{HTML}{A52D39}
\definecolor{reperereponse}{HTML}{14828B}
\renewcommand{\arraystretch}{1.3}
\begin{tabularx}{\textwidth}{@{}*{3}{>{\raggedright\arraybackslash}X}@{}}
\toprule
\textbf{Question à résoudre} & \textcolor{reperecritique}{\textbf{Méthode et apport}} & \textcolor{reperereponse}{\textbf{Résultat à examiner et limite}} \\
\midrule
Comment représenter l'effet d'un état de conformité que le forfait ne distingue pas ?
& Relier un état supposé aux paramètres d'une même chaîne de pertes (chapitres~\ref{chap:socle} à~\ref{chap:multietats}).
& Comparer les capitaux annuels (\S\ref{sec:chiffre-de-tete}) ; les écarts entre états sont posés, sans effet causal établi. \\
\addlinespace
Que reste-t-il du calcul si les données ne donnent pas le sens des transmissions ?
& Examiner les sources, séparer force et asymétrie, puis explorer les directions admissibles (chapitres~\ref{ch:identifiabilite} et~\ref{chap:partielle}).
& Une plage aux sommets, à force commune fixée ; ni intervalle de confiance ni borne globale démontrée. \\
\addlinespace
Quelle correction comparer, et quelles données manquent pour décider ?
& Calculer les effets de fermeture, puis éprouver les coûts et la calibration (chapitres~\ref{chap:resultats} et~\ref{ch:robustesse}).
& Gains conditionnels et besoins de collecte ; l'efficacité réelle d'une remédiation reste à mesurer. \\
\bottomrule
\end{tabularx}
\endgroup
\caption{Les questions qui guident le mémoire et les réponses à vérifier. Les renvois indiquent
où la méthode est construite puis éprouvée ; la dernière colonne en délimite la portée.}
\label{tab:apports-institut}
\end{table}


La première étape consiste donc à déterminer les entrées réellement observées. Le chapitre
suivant distingue les pertes disponibles des états et des transmissions qui resteront posés.

\nocite{BohmeKataria2006, HerathHerath2011, HillairetLopez2021,
  HillairetLopezDOultremontSpoorenberg2022, XuHua2017, Baldwin2017,
  BessyRolandBoumezouedHillairet2021, BoumezouedCherkaouiHillairet2023,
  HillairetReveillacRosenbaum2021, Peng2016, Edwards2016, LuZhangZhu2024,
  FarkasLopezThomas2021, Vasicek2002, Leontief1936}
